\documentclass[12pt]{article}
\usepackage[utf8]{inputenc}
\usepackage[spanish]{babel}
\usepackage{amsmath, amssymb}
\usepackage{geometry}
\geometry{a4paper, margin=2.5cm}
\usepackage{xcolor}
\usepackage{tcolorbox}
\tcbuselibrary{skins}

\definecolor{azulIPn}{HTML}{003DA5}
\definecolor{verdeMatem}{HTML}{1de9b6}
\definecolor{rojoSuper}{HTML}{ff6b6b}

\newtcolorbox{definicion}[1]{
  colback=azulIPn!10, colframe=azulIPn,
  fonttitle=\bfseries, title=Definición: #1
}
\newtcolorbox{nota}{
  colback=rojoSuper!10, colframe=rojoSuper!60!black,
  fonttitle=\bfseries, title=Nota importante
}

\title{\textcolor{azulIPn}{\textbf{Nota 02 — Vectores y Operaciones Básicas}} \\
       \large Introducción al Cálculo Vectorial}
\author{Daniel Tovar Abraham \\ ESCOM -- IPN}
\date{Servicio Social 2026}

\begin{document}
\maketitle

% -----------------------------------------------
\section{Vectores en $\mathbb{R}^3$}
% -----------------------------------------------

Un vector en $\mathbb{R}^3$ se escribe como:
\[
  \vec{v} = (v_1,\, v_2,\, v_3) = v_1\,\hat{i} + v_2\,\hat{j} + v_3\,\hat{k}
\]

donde $\hat{i}$, $\hat{j}$, $\hat{k}$ son los vectores unitarios en las 
direcciones $x$, $y$, $z$ respectivamente.

% -----------------------------------------------
\section{Operaciones básicas}
% -----------------------------------------------

\subsection{Suma y resta}
\[
  \vec{u} \pm \vec{v} = (u_1 \pm v_1,\; u_2 \pm v_2,\; u_3 \pm v_3)
\]

\subsection{Magnitud (norma)}
\[
  \|\vec{v}\| = \sqrt{v_1^2 + v_2^2 + v_3^2}
\]

\subsection{Vector unitario}
\[
  \hat{v} = \frac{\vec{v}}{\|\vec{v}\|}
\]

\begin{nota}
  Un vector unitario tiene magnitud 1. Se usa para indicar 
  \textbf{solo la dirección}, sin importar la magnitud.
  En las animaciones: $\|\hat{n}\| = 1$ es el vector normal unitario 
  a una superficie.
\end{nota}

% -----------------------------------------------
\section{Producto punto (escalar)}
% -----------------------------------------------

\begin{definicion}{Producto punto}
  \[
    \vec{u} \cdot \vec{v} = u_1 v_1 + u_2 v_2 + u_3 v_3 
                          = \|\vec{u}\|\,\|\vec{v}\|\cos\theta
  \]
  donde $\theta$ es el ángulo entre los vectores.
\end{definicion}

\textbf{Propiedades importantes:}
\begin{itemize}
  \item Si $\vec{u} \cdot \vec{v} = 0$ entonces $\vec{u} \perp \vec{v}$
  \item $\vec{u} \cdot \vec{u} = \|\vec{u}\|^2$
  \item Conmutativo: $\vec{u} \cdot \vec{v} = \vec{v} \cdot \vec{u}$
\end{itemize}

% -----------------------------------------------
\section{Producto cruz (vectorial)}
% -----------------------------------------------

\begin{definicion}{Producto cruz}
  \[
    \vec{u} \times \vec{v} = 
    \begin{vmatrix}
      \hat{i} & \hat{j} & \hat{k} \\
      u_1 & u_2 & u_3 \\
      v_1 & v_2 & v_3
    \end{vmatrix}
    = (u_2 v_3 - u_3 v_2)\,\hat{i} 
    - (u_1 v_3 - u_3 v_1)\,\hat{j} 
    + (u_1 v_2 - u_2 v_1)\,\hat{k}
  \]
\end{definicion}

\textbf{Propiedades importantes:}
\begin{itemize}
  \item El resultado es un vector \textbf{perpendicular} a ambos
  \item $\|\vec{u} \times \vec{v}\| = \|\vec{u}\|\,\|\vec{v}\|\sin\theta$
  \item \textbf{No} es conmutativo: $\vec{u} \times \vec{v} = -(\vec{v} \times \vec{u})$
\end{itemize}

\bigskip
\begin{nota}
  El producto cruz se usa para calcular el \textbf{vector normal} a una 
  superficie paramétrica:
  \[
    \hat{n} = \frac{\vec{r}_u \times \vec{r}_v}{\|\vec{r}_u \times \vec{r}_v\|}
  \]
  Esto es exactamente lo que aparece en la animación de la Cinta de Möbius.
\end{nota}

\end{document}